\emph{Preliminaries.} The blades are a basis of $\Cl_{k}(\ell,\ell)$ and the adjoint $\dagger$, the linear extension of the inverse on blades, is an anti-automorphism that fixes each $e_{j}$ and negates each $f_{j}$ (Theorem \claimref{synthesis.signed-directions-and-the-integer-clifford-algebra}); so each $p_{j} = f_{j}e_{j}$ is self-adjoint of square $1$, the $p_{j}$ commute with each other, and $p_{j}$ anticommutes with $e_{j}$, $f_{j}$ and commutes with the other generators. Hence $\gamma = p_{1}\cdots p_{m}$ is self-adjoint of square $1$, anticommutes with the generators of the first $m$ modes and commutes with the others. The twisted generators are those of Proposition \claimref{tower.commuting-copies}, with the modes of party $2$ as the ground algebra and the $m$ modes of party $1$ as the adjoined planes (its product of plane pseudoscalars is $\pm\gamma$); the three facts needed are proved here directly, over any $k$. (1) A twisted generator $\gamma g$ ($g = e_{j}$ or $f_{j}$, $j > m$) commutes with each generator $h$ of party $1$: $h\gamma g = -\gamma hg = \gamma gh$. So $\mathcal{A}$ and $\mathcal{B}$ commute elementwise. (2) $(\gamma g)(\gamma g') = \gamma^{2}gg' = gg'$, so the twisted generators satisfy the relations of the generators of $\Cl_{k}(n,n)$, and an ordered product of $r$ of them is $\gamma^{r}$ times the corresponding blade $g_{R}$ of party $2$'s generators, with $\gamma^{r} \in \{1, \gamma\}$. (3) $\gamma$ is $\pm$ the blade of all $2m$ generators of party $1$, and a blade of $\mathcal{A}$ times $\gamma$ is $\pm$ the complementary blade of $\mathcal{A}$. So the $4^{m}\cdot 4^{n}$ products of a blade $a_{T}$ of $\mathcal{A}$ with an ordered product of twisted generators are $\pm a_{T}g_{R}$ for $|R|$ even and $\pm a_{T^{c}}g_{R}$ for $|R|$ odd: signed blades of $\Cl_{k}(\ell,\ell)$, each blade arising once. Hence multiplication $\mathcal{A} \otimes \mathcal{B} \to \Cl_{k}(\ell,\ell)$ is a linear bijection, $\mathcal{A} \cong \Cl_{k}(m,m)$ and $\mathcal{B} \cong \Cl_{k}(n,n)$ (the check verifies the count for $\ell = 4$, $m = 2$, in the form of $256$ monomials, and for $\ell = 2$). $\mathcal{A}$ is closed under $\dagger$, and so is $\mathcal{B}$: $(\gamma e_{j})^{\dagger} = e_{j}\gamma = \gamma e_{j}$ and $(\gamma f_{j})^{\dagger} = -f_{j}\gamma = -\gamma f_{j}$. Restricted to $\mathcal{B}$, $\dagger$ is the anti-automorphism that fixes the twisted generators of square $1$ and negates those of square $-1$.

\emph{(i).} $S$ and $D$ commute with $B_{0}$, $B_{1}$, and $F_{+} + F_{-} = 1$, $F_{+} - F_{-} = B_{0}$, $G_{+} + G_{-} = 1$, $G_{+} - G_{-} = B_{1}$. So $(S - t)^{2}F_{+} + (S + t)^{2}F_{-} = (S^{2} + t^{2})(F_{+} + F_{-}) - 2tS(F_{+} - F_{-}) = S^{2} + t^{2} - 2tSB_{0}$, and likewise $(D - t)^{2}G_{+} + (D + t)^{2}G_{-} = D^{2} + t^{2} - 2tDB_{1}$. Since $A_{0}^{2} = A_{1}^{2} = 1$, $S^{2} + D^{2} = 2A_{0}^{2} + 2A_{1}^{2} = 4$, and $SB_{0} + DB_{1} = C$; the sum is $4 + 2t^{2} - 2tC$. Each $F$ is a self-adjoint idempotent ($B_{j}^{2} = 1$, $B_{j}^{\dagger} = B_{j}$), each $s = S \mp t$ or $D \mp t$ is self-adjoint and commutes with $F$, so with $y = sF$, $yy^{\dagger} = sFFs = s^{2}F$. Nothing else was used.

\emph{(ii).} Apply $\lambda$ to (i): under (CP), $4 + 2t^{2} - 2t\lambda(C) \ge 0$ for every $t \in k$. At $t = \tfrac12\lambda(C)$ this is $4 - \tfrac12\lambda(C)^{2} \ge 0$, so $\lambda(C)^{2} \le 8$; over $\Q$ equality is excluded because $8$ is not the square of a rational. Under (P) each $\lambda(s^{2}F) = \lambda(yy^{\dagger}) \ge 0$, so (P) gives (CP). If $a \mapsto \lambda(aF)$ is non-negative on the elements $aa^{\dagger}$ of $\mathcal{A}$, then $\lambda(s^{2}F) = \lambda(ss^{\dagger}F) \ge 0$. For one mode per party, $k$ a subfield of $\R$ and directions $A_{0}, A_{1}, B_{0}, B_{1}$ (notation of the cited Theorem): a self-adjoint $s \in \mathrm{span}\{1, A_{0}, A_{1}\}$ is $\alpha + \xi X_{1} + \zeta Z_{1}$, and $s^{2} = \alpha^{2} + \nu^{2} + 2\alpha(\xi X_{1} + \zeta Z_{1})$ with $\nu^{2} = \xi^{2} + \zeta^{2}$. Put $w = \lambda(F)$ and let $g$ be the linear form $(\xi, \zeta) \mapsto \lambda((\xi X_{1} + \zeta Z_{1})F)$ on $k^{2}$. For a direction $d = xX_{1} + zZ_{1}$, (L) at the outcomes $(\pm d, F)$ gives $w \pm g(x,z) \ge 0$; the points of the unit circle with coordinates in $k$ are dense in it, so $|g(\xi,\zeta)| \le w\nu$ for all $(\xi,\zeta)$, and $\lambda(s^{2}F) = (\alpha^{2} + \nu^{2})w + 2\alpha\,g(\xi,\zeta) \ge (|\alpha| - \nu)^{2}w \ge 0$. The box functional: with $A_{0} = X_{1}$, $A_{1} = Z_{1}$ anticommuting, $S^{2} = 2$ and $(S - 1)^{2}F_{+} = \tfrac12(3 - 2S)(1 + B_{0})$; its scalar part is $\tfrac32$ and $\langle C_{0}(S - 1)^{2}F_{+}\rangle_{0} = -\langle SB_{0}SB_{0}\rangle_{0} = -\langle S^{2}\rangle_{0} = -2$ (the other products have no scalar part), so the value is $-\tfrac12$; the other three terms are alike, the check evaluates all four, and their sum is $4 + 2 - 2\cdot 4 = -2$ as (i) requires for a functional of value $4$.

\emph{(iii).} $\mathcal{A}$ with $\dagger$ is the algebra of $m \ge 2$ modes (the blades in $e_{j}, f_{j}$, $j \le m$, are a basis of $\mathcal{A}$ by Theorem \claimref{synthesis.signed-directions-and-the-integer-clifford-algebra}, and $\dagger$ restricts). By clause (i) of Theorem \claimref{synthesis.additive-weights-on-the-rational-projectors-are-linear-in-every-size-from-four} there is an isomorphism $\varphi$ of $\mathcal{A}$ onto $\Mat_{N}(\Q)$, $N = 2^{m} \ge 4$, with $\varphi(x^{\dagger}) = \varphi(x)^{T}$, under which the self-adjoint idempotents of $\mathcal{A}$ are the symmetric idempotent matrices; write $\Pi_{r}$ for the element with matrix $rr^{T}/(r \cdot r)$, $r \in \Q^{N}$ nonzero. For fixed $F$ the function $E \mapsto \mu(EF)$ on the self-adjoint idempotents of $\mathcal{A}$ adds on orthogonal pairs, so by clause (iii) of that Theorem (for $m = 2$, clause (iii) of Theorem \claimref{synthesis.additive-weights-on-the-rational-projectors-of-two-modes-are-linear}) it is the restriction of a $\Q$-linear $L_{F}(x) = \sum_{ij}\rho^{F}_{ij}\varphi(x)_{ij}$ on $\mathcal{A}$ with $\rho^{F}_{ij} = \rho^{F}_{ji}$, so that $L_{F}(x^{\dagger}) = L_{F}(x)$. For self-adjoint $s \in \mathcal{A}$ the matrix of $s^{2} = ss^{\dagger}$ is $\sum_{j}c_{j}c_{j}^{T}$ over the nonzero columns $c_{j}$ of $\varphi(s)$, so $s^{2} = \sum_{j}(c_{j} \cdot c_{j})\Pi_{c_{j}}$ and $L_{F}(s^{2}) = \sum_{j}(c_{j} \cdot c_{j})\,\mu(\Pi_{c_{j}}F) \ge 0$. The self-adjoint idempotents span the self-adjoint elements of $\mathcal{A}$ (every symmetric rational matrix is a rational combination of matrices $rr^{T}$), so the hypothesis $\mu(EF_{+}) + \mu(EF_{-}) = \mu(E)$ gives $L_{F_{+}} + L_{F_{-}} = L_{1}$ on self-adjoint elements, and likewise for $G_{\pm}$; $L_{1}(1) = \mu(1) = 1$. With $E_{i,\pm} = \tfrac12(1 \pm A_{i})$, $A_{i} = E_{i,+} - E_{i,-}$ and $(L_{F_{+}} - L_{F_{-}})(A_{i}) = \sum_{a,b}ab\,\mu(E_{i,a}F_{b}) = c_{i0}$, and likewise $(L_{G_{+}} - L_{G_{-}})(A_{i}) = c_{i1}$. Therefore
\[
0 \le L_{F_{+}}\bigl((S - t)^{2}\bigr) + L_{F_{-}}\bigl((S + t)^{2}\bigr) + L_{G_{+}}\bigl((D - t)^{2}\bigr) + L_{G_{-}}\bigl((D + t)^{2}\bigr) = L_{1}(S^{2} + D^{2}) + 2t^{2} - 2t\bigl(c_{00} + c_{10} + c_{01} - c_{11}\bigr) = 4 + 2t^{2} - 2tV ,
\]
and $t = \tfrac12 V$ gives $V^{2} \le 8$, hence $V^{2} < 8$ over $\Q$. The check exhibits one $\mu$ satisfying the hypotheses on three modes ($m = 2$), coming from a functional that is not a state, with $V = -\tfrac{14}{5}$, and one on four modes ($m = 3$), with $V = -\tfrac{34}{13}$, and verifies the bound and the sum of the four values for them; the passage from additive to linear is the stored theorems' and is not recomputed.

\emph{(iv).} By the Preliminaries the bilinear map $(a, b) \mapsto ab^{\dagger}$ on $\mathcal{A} \times \mathcal{B}$ factors through one linear map $\theta$ of $\Cl_{k}(\ell,\ell)$. $\theta^{2}(ab) = ab^{\dagger\dagger} = ab$ and $\theta(1) = 1$. For $x = ab$: $x^{\dagger} = b^{\dagger}a^{\dagger} = a^{\dagger}b^{\dagger}$, so $\theta(x^{\dagger}) = a^{\dagger}b = (ab^{\dagger})^{\dagger} = \theta(x)^{\dagger}$. For $a', a'' \in \mathcal{A}$: $\theta(a'aba'') = \theta((a'aa'')b) = a'aa''b^{\dagger} = a'\theta(ab)a''$. For $b' \in \mathcal{B}$ self-adjoint: $\theta(b'abb') = \theta(a\,b'bb') = a(b'bb')^{\dagger} = a\,b'b^{\dagger}b' = b'\theta(ab)b'$. The parity automorphism preserves $\mathcal{A}$ and $\mathcal{B}$ (it negates their generators) and commutes with $\dagger$, hence with $\theta$. Linearity extends each identity from products to all $x$. For $\lambda^{\theta} = \lambda \circ \theta$: it is linear and $\lambda^{\theta}(1) = \lambda(1)$; if $\lambda$ is real, $\lambda^{\theta}(x^{\dagger}) = \lambda(\theta(x)^{\dagger}) = \lambda(\theta(x))$; if $\lambda$ is even it vanishes on $\theta$ of an odd element, which is odd. For self-adjoint $b$, $\theta(ab) = ab$, so $\lambda^{\theta}(ab) = \lambda(ab)$ for every $a \in \mathcal{A}$. For a product of two self-adjoint elements $b, b'$ of $\mathcal{B}$ that do not commute this fails: $\theta(abb') = ab'b \ne abb'$ in general, and in (v) $(X_{1}Z_{1})(X_{2}Z_{2}) = -K$ is negated by $\theta$, with $\lambda_{-1}(-K) = 1$ and $\lambda_{1}(-K) = -1$ (check). In a sequence $P_{1}\cdots P_{r-1}P_{r}P_{r-1}\cdots P_{1}$ of self-adjoint idempotents from $\mathcal{A}$ or $\mathcal{B}$ the factors from $\mathcal{A}$ commute with those from $\mathcal{B}$, so the product is $ab$ with $a$ the product of the factors from $\mathcal{A}$ in their order and $b$ that of the factors from $\mathcal{B}$; each is a palindrome of self-adjoint elements, hence self-adjoint, and the two functionals agree on it. The elements $s^{2}F$ of (CP) are of the form $ab$ with $b = F$ self-adjoint. For a local observable $o$ of square $1$ the identities above give $\theta(oxo) = o\,\theta(x)\,o$, so $\theta \circ E_{o} = E_{o} \circ \theta$.

\emph{(v).} $X_{1}, Z_{1}$ generate $\mathcal{A}$ and $X_{2} = e_{2}p_{1} = \gamma e_{2}$, $Z_{2} = p_{2} = (\gamma f_{2})(\gamma e_{2})$ lie in $\mathcal{B}$, so $U$ and $V$ are products of an element of $\mathcal{A}$ and a self-adjoint element of $\mathcal{B}$ and are fixed by $\theta$, while $Y_{2}^{\dagger} = Z_{2}X_{2} = -Y_{2}$ gives $\theta(K) = -Y_{1}Y_{2}^{\dagger} = -K$. The sixteen monomials in $X_{j}, Z_{j}, Y_{j}$ are signed blades, pairwise orthogonal for $\langle xy\rangle_{0}$ up to sign, and $\theta$ multiplies each by $\pm 1$; so $\langle M\,\theta(x)\rangle_{0} = \langle\theta(M)x\rangle_{0}$ and $\lambda_{c}^{\theta} = \lambda_{-c}$. The weights: $U$ and $V$ commute, have square $1$ and $UV = -K$ (check), so $U(1 + uU) = u(1 + uU)$, $V(1 + vV) = v(1 + vV)$ and $(1 - U - V + cK)\,\tfrac14(1 + uU)(1 + vV) = (1 - u - v - c\,uv)\,\tfrac14(1 + uU)(1 + vV)$; the idempotent has scalar part $\tfrac14$, so $\lambda_{c}$ gives it $\tfrac14(1 - u - v - c\,uv)$: $-\tfrac14(1 + c)$ at $(+,+)$ and $\tfrac14(1 + c)$ at $(+,-)$, one of which is negative on a self-adjoint idempotent unless $c = -1$. At $c = -1$ the density is $(1 - U)(1 - V) = 4P_{s}$ with $P_{s}$ a self-adjoint idempotent, and $\lambda_{-1}(xx^{\dagger}) = 4\langle P_{s}xx^{\dagger}P_{s}\rangle_{0} = 4\langle(P_{s}x)(P_{s}x)^{\dagger}\rangle_{0} \ge 0$. So $\lambda_{c}$ is non-negative on all $xx^{\dagger}$ exactly for $c = -1$ (this is also clause (iv) of the cited Theorem; it is proved here and that result is not used). At $c = 1$ the weights are $-\tfrac12, \tfrac12, \tfrac12, \tfrac12$; the inertia is computed by the check. Two copies. The elements $X_{j}, Z_{j}$, $j = 1, \dots, 4$, generate four commuting copies of the two-by-two matrices, the $256$ monomials are signed blades and each blade arises once (check), $\mathcal{A}$ is generated by those with $j = 1, 2$ and $\mathcal{B}$ by those with $j = 3, 4$ ($X_{3} = \pm\gamma e_{3}$, $Z_{3} = \pm(\gamma f_{3})(\gamma e_{3})$, $X_{4} = \pm(\gamma e_{4})Z_{3}$; check). The scalar part of a product of a monomial in the modes $1, 3$ and a monomial in the modes $2, 4$ is the product of the scalar parts, so $x \mapsto \langle M^{(13)}_{c}M^{(24)}_{c}x\rangle_{0}$ is the product of the two copies of $\lambda_{c}$, with $\lambda_{c}(X_{a}X_{b}) = \lambda_{c}(Z_{a}Z_{b}) = -1$, $\lambda_{c}(Y_{a}Y_{b}) = -c$ and $0$ on the other non-scalar monomials of a pair (check, all $256$ monomials). Expanding the two idempotents, the terms with non-zero value are $1$; $ss'(X_{1}X_{2})(X_{3}X_{4}) = ss'(X_{1}X_{3})(X_{2}X_{4})$, of value $ss'$; $tt'(Z_{1}Z_{3})(Z_{2}Z_{4})$, of value $tt'$; and $ss'tt'(X_{1}X_{2}Z_{1}Z_{2})(X_{3}X_{4}Z_{3}Z_{4}) = ss'tt'(Y_{1}Y_{3})(Y_{2}Y_{4})$, of value $ss'tt'c^{2}$; a term such as $(X_{1}X_{2})(Z_{3}Z_{4}) = (X_{1}Z_{3})(X_{2}Z_{4})$ has value $0$. This gives the formula, verified by the check for all sixteen sign choices at four values of $c$. The reversal on party $2$ changes the sign of $Y_{3}$ and $Y_{4}$, so it sends the product for $c$ to that for $-c$; for $c = -1$ the density $M^{(13)}_{-1}M^{(24)}_{-1}$ is $16$ times a self-adjoint idempotent $\Pi$ (check), so the functional is $x \mapsto 16\langle\Pi x\rangle_{0}$, non-negative on $xx^{\dagger}$ because $\langle\Pi xx^{\dagger}\rangle_{0} = \langle(\Pi x)(\Pi x)^{\dagger}\rangle_{0}$ is a sum of squares of coefficients; and (iv) gives the last statement.

\emph{(vi).} By clause (i) of Theorem \claimref{synthesis.additive-weights-on-the-rational-projectors-of-two-modes-are-linear} identify $\Cl_{\Q}(2,2)$ with $\Mat_{4}(\Q)$, $\dagger$ with the transpose and the atoms with the matrices $\Pi_{r} = rr^{T}/(r \cdot r)$. Write $\lambda(x) = \sum_{ij}\rho_{ij}m(x)_{ij}$ and let $\sigma$ be the symmetric part of $\rho$. Then $\lambda(xx^{\dagger}) = \sum_{j}c_{j}^{T}\sigma c_{j}$ over the columns of $m(x)$ and $\lambda(\Pi_{r}) = r^{T}\sigma r/(r \cdot r)$, so $\lambda$ is non-negative on all $xx^{\dagger}$ exactly when $r^{T}\sigma r \ge 0$ for all $r \in \Q^{4}$. The elements with $uu^{\dagger} = 1$ are the rational orthogonal matrices and $u\Pi_{r}u^{\dagger} = \Pi_{ur}$. Let $P = \Pi_{r_{0}}$. If $r \cdot r = q^{2}(r_{0} \cdot r_{0})$ with $q$ rational, the reflection $1 - 2dd^{T}/(d \cdot d)$, $d = qr_{0} - r$ (or the identity if $d = 0$), is orthogonal and carries $qr_{0}$ to $r$, so $\Pi_{r} = uPu^{\dagger}$. Hence the hypothesis gives $r^{T}\sigma r \ge 0$ on the set $N$ of vectors whose norm is $r_{0} \cdot r_{0}$ times a rational square. $N$ is dense in $\Q^{4}$: read vectors as rotors $x = x_{0} + x_{1}B_{1} + x_{2}B_{2} + x_{3}B_{3}$ of the even subalgebra of $\Cl_{\Q}(3,0)$ (\claimref{std-algebra.quaternions}), whose norm $x\tilde{x}$ is multiplicative; the rotors of norm a rational square are the rational multiples of the rational points of the unit sphere, which are dense in it (stereographic parametrisation), so they are dense in $\Q^{4}$, and left multiplication by the rotor of $r_{0}$, a linear bijection that multiplies norms by $r_{0} \cdot r_{0}$, carries them onto $N$. If $r^{T}\sigma r < 0$ for some rational $r$, the same holds on a neighbourhood of $r$, which meets $N$; so $r^{T}\sigma r \ge 0$ for all $r$. The converse is immediate, $uPu^{\dagger} = (uP)(uP)^{\dagger}$. The instance is a computation: $u^{\dagger} = u$ since $1, Z_{2}, X_{1}, X_{1}Z_{2}$ are self-adjoint ($X_{1}$ and $Z_{2}$ commute), $u^{2} = 1$ by expansion, $uZ_{1}u = Z_{1}Z_{2}$ and $uX_{2}u = X_{1}X_{2}$ (check, in the form $u\,\tfrac14(1 + Z_{1})(1 + X_{2})\,u = \tfrac14(1 + U)(1 + V)$), and the weights follow from (v).

\emph{(vii).} $I_{4}$ is a product of four generators that anticommute with each mode vector, so it commutes with the mode vectors; $I_{4}^{2} = -1$ by the signature, and the blade inverse gives $I_{4}^{\dagger} = I_{4}^{-1} = -I_{4}$. $Y_{j}^{\dagger} = -Y_{j}$ and $Y_{j}^{2} = -1$, so $(I_{4}Y_{j})^{\dagger} = Y_{j}^{\dagger}I_{4}^{\dagger} = Y_{j}I_{4} = I_{4}Y_{j}$ and $(I_{4}Y_{j})^{2} = I_{4}^{2}Y_{j}^{2} = 1$. $Y_{1}$ and $Y_{2}$ commute, so $(I_{4}Y_{1})(I_{4}Y_{2}) = -Y_{1}Y_{2} = K$. The check verifies each relation in $\Cl_{\Q}(3,5)$.
